\clearpage
\begingroup
\fancyhead[L]{\small BEMERKUNGEN}
\renewcommand{\headrulewidth}{0pt}
\begingroup
\newcommand{\BVIJ}{\text{\mdseries\AstroSymbols ♃}}
\newcommand{\BVIP}{\text{\mdseries\AstroSymbols ⚴}}
\newcommand{\BVIsin}{\mathop{\mathit{sin}}}
\par\noindent\hspace*{\parindent}Wirklich vollendet wurde die ganze Rechnung im Juli 1813; s. Brief an
\textsc{Olbers} S. 422.

\par\noindent\hspace*{\parindent}Die Seite 529 abgedruckten Elementensysteme, sowie die Berechnung der
Lage der Bahnen gegen einander stehen im Handbuch Bc, ebenso wie die »Präparation
der Jupiters- und Pallasörter«, d. h. die ebenfalls S. 529 erwähnte Berechnung
der Coordinaten \(r\), \(v\); \(r'\), \(w'\), \(\beta'\), \(v'\).
Die weitere Rechnung ist vollständig auf einer grossen Zahl einzelner Blätter
ausgeführt, denen die Zahlen der Tafeln Seite 530--541, sowie die der
Rechnungsbeispiele für das Argument \(5M'-2M\), Seite 542, 545, 548, 550, 552
entnommen sind. Die schliesslichen Resultate für die Störungen, Seite 543--545,
546--547, 548--549, 550--552 (nur \(\delta\log a\)), 553--554, 555--556,
stehen wieder in dem erwähnten Handbuch; nur die Coefficienten von
\(\delta\int n\,\mathrm{d}t\) haben aus Blättern entnommen werden müssen.
Bei dem Ausdruck für \(\delta\tilde\omega\), Seite 547, sind leider aus Versehen
einige Glieder vergessen worden, die darum hier noch angeführt werden sollen, nemlich:
\[
\begin{array}{r@{\ }r@{\ }r}
+\mathit{1\overset{\prime\prime}{,}24}\,\BVIsin(\mathit{14}\BVIJ-\mathit{5}\BVIP+\mathit{253^0}&\mathit{0'}&\mathit{39''})\\
+\mathit{0{,}24}\,\BVIsin(\mathit{14}\BVIJ-\mathit{4}\BVIP+\mathit{251}&\mathit{4}&\mathit{55})\\
+\mathit{0{,}13}\,\BVIsin(\mathit{15}\BVIJ-\mathit{10}\BVIP+\mathit{116}&\mathit{58}&\mathit{57})\\
+\mathit{0{,}19}\,\BVIsin(\mathit{15}\BVIJ-\mathit{9}\BVIP+\mathit{106}&\mathit{32}&\mathit{59})\\
+\mathit{0{,}29}\,\BVIsin(\mathit{15}\BVIJ-\mathit{8}\BVIP+\mathit{99}&\mathit{41}&\mathit{37})\\
+\mathit{0{,}47}\,\BVIsin(\mathit{15}\BVIJ-\mathit{7}\BVIP+\mathit{95}&\mathit{9}&\mathit{38})\\
+\mathit{2{,}59}\,\BVIsin(\mathit{15}\BVIJ-\mathit{6}\BVIP+\mathit{91}&\mathit{52}&\mathit{25})\\
+\mathit{0{,}35}\,\BVIsin(\mathit{15}\BVIJ-\mathit{5}\BVIP+\mathit{268}&\mathit{44}&\mathit{35})\\
+\mathit{0{,}10}\,\sin(\mathit{16}\BVIJ-\mathit{10}\BVIP+\mathit{125}&\mathit{50}&\mathit{29})\\
+\mathit{0{,}15}\,\BVIsin(\mathit{16}\BVIJ-\mathit{9}\BVIP+\mathit{117}&\mathit{59}&\mathit{33})\\
+\mathit{0{,}22}\,\BVIsin(\mathit{16}\BVIJ-\mathit{8}\BVIP+\mathit{109}&\mathit{38}&\mathit{55})\\
+\mathit{0{,}42}\,\BVIsin(\mathit{16}\BVIJ-\mathit{7}\BVIP+\mathit{107}&\mathit{50}&\mathit{23})
\end{array}
\]
\[
\begin{array}{r@{\ }r@{\ }r}
+\mathit{1\overset{\prime\prime}{,}10}\,\BVIsin(\mathit{16}\BVIJ-\mathit{6}\BVIP+\mathit{283^0}&\mathit{10'}&\mathit{1''})\\
+\mathit{0{,}13}\,\BVIsin(\mathit{16}\BVIJ-\mathit{5}\BVIP+\mathit{281}&\mathit{17}&\mathit{22})\\
+\mathit{0{,}12}\,\BVIsin(\mathit{17}\BVIJ-\mathit{9}\BVIP+\mathit{129}&\mathit{36}&\mathit{2})\\
+\mathit{0{,}20}\,\BVIsin(\mathit{17}\BVIJ-\mathit{8}\BVIP+\mathit{125}&\mathit{14}&\mathit{12})\\
+\mathit{0{,}50}\,\BVIsin(\mathit{17}\BVIJ-\mathit{7}\BVIP+\mathit{120}&\mathit{22}&\mathit{4})\\
+\mathit{0{,}22}\,\BVIsin(\mathit{17}\BVIJ-\mathit{6}\BVIP+\mathit{297}&\mathit{40}&\mathit{56})\\
+\mathit{0{,}14}\,\BVIsin(\mathit{18}\BVIJ-\mathit{8}\BVIP+\mathit{137}&\mathit{75}&\mathit{16})\\
+\mathit{0{,}14}\,\BVIsin(\mathit{19}\BVIJ-\mathit{8}\BVIP+\mathit{149}&\mathit{16}&\mathit{20})\\
+\mathit{0{,}16}\,\BVIsin(\mathit{19}\BVIJ-\mathit{7}\BVIP+\mathit{146}&\mathit{41}&\mathit{38})\\
+\mathit{0{,}22}\,\BVIsin(\mathit{20}\BVIJ-\mathit{8}\BVIP+\mathit{161}&\mathit{17}&\mathit{24})\\
+\mathit{0{,}16}\,\BVIsin(\mathit{21}\BVIJ-\mathit{8}\BVIP+\mathit{173}&\mathit{18}&\mathit{28}).
\end{array}
\]

\par\noindent\hspace*{\parindent}Die numerischen Werthe von Notiz [9.], S. 557, stehen auf den einzelnen
Blättern zwischen den Integrationen für die übrigen Gleichungen; zum Theil
mussten sie ergänzt werden.

\par\noindent\hspace*{\parindent}Besonderes Interesse wird die Notiz [10.] beanspruchen; sie ist aus
einigen ganz kurzen Notizen auf zwei unscheinbaren Zettelchen zusammengestellt,
von denen der eine aus einem Convolut von losen Zetteln
\SourceLeaf{614}stammt, die nach \textsc{Gauss}' Tode sich in seinem Papierkorb vorgefunden haben sollen.
Wenn auch das Aussehen dieses Zettels eine derartige Herkunft wahrscheinlich
macht, so ist doch mit Rücksicht auf den Inhalt kaum anzunehmen, dass er kurz
vor \textsc{Gauss}' Tode in den Papierkorb geworfen wurde; bei den wiederholten
Durchsichten des Nachlasses hat es sich naturgemäss nicht ganz vermeiden lassen,
dass die ursprüngliche Ordnung der einzelnen aufgefundenen Zettel verändert wurde.
Aus welcher Zeit die beiden Zettel stammen, dürfte kaum sicher festzustellen
sein; nach Ausweis des vorstehend abgedruckten »Journal über die Rechnungen an
den Pallasstörungen« konnte \textsc{Gauss} erst in der Zeit zwischen dem 20. Juli
und 10. August 1812 zur Kenntnis des numerischen Werths des Coefficienten
\(\alpha\) (S. 558) gelangt sein, also nach Verfassung des Briefes an
\textsc{Bessel}, S. 422, aber vor Abfassung des Passus in der »Exposition«, S. 472.
Für die Annahme einer spätern Entstehungszeit beider Zettel spricht vielleicht
der Umstand, dass \textsc{Gauss} bei Entnahme der numerischen Werthe aus seinen
frühern Rechnungen ein Versehen begangen hat, indem er die Constante \(A\)
unrichtig ansetzte, nemlich
\[
A=198^0\,39'\,8'',
\]
\noindent womit er findet
\[
\begin{aligned}
U&=324^0\,15'\,57''\\
\log C&=4{,}26300
\end{aligned}
\]
\noindent und für die Periode
\[
1026{,}17\ \text{Pallasumläufe}=399{,}07\ \text{Jupitersumläufe}.
\]

\par\noindent\hspace*{\parindent}Die Grenzen, zwischen denen \(n\) und \(u\) oscilliren (S. 559), hat
\textsc{Gauss} nicht berechnet, mit dem obigen Werthe von \(A\) würde folgen:
\[
\begin{array}{ccccc@{\ }l}
\text{für}&\text{die}&\text{Grenzen}&\text{von}&n:&\dfrac{18}{7}n'\pm0\overset{\prime\prime}{,}07219\\[3pt]
\text{»}&\text{»}&\text{»}&\text{»}&u:&\pm37^0\,35'.
\end{array}
\]
\noindent Beim Abdruck sind die verbesserten Werthe eingesetzt worden. Der Unterschied
ist recht beträchtlich; man muss aber bedenken, dass das ganze Resultat sehr
precär ist, weil der Wert von \(\mu\) und damit der von \(C\) sehr unsicher zu bestimmen ist.

\par\noindent\hspace*{\parindent}In den Originalen befinden sich nur die Gleichungen Seite 558, Zeile 4,
6, 14 von oben und 8, 6, 5, 4, 3 von unten, Seite 559, Zeile 5, sowie die numerischen
Werthe von \(\alpha\), \(A\), \(18\BVIJ-7\BVIP\), \(U\), \(14\alpha\), \(\mu\)
und die Notiz »\textit{1026,17 Pallasumläufe, 399,07 Jupitersumläufe}«.

\par\noindent\hspace*{\parindent}Die Untersuchungen über die Berechnung der Argumente mit Berücksichtigung
der grossen Gleichung, Notiz [11.], finden sich grösstentheils auf Blättern,
sonst im Handbuch Bc; den Blättern sind jedoch nur die Ausdrücke \(\mathrm{f}(u)\)
und \(\mathrm{F}(u)\), Seite 559, und dem Handbuche die beiden Relationen 1) und 3)
entnommen. Die Bearbeitung musste auch hier sehr frei gehalten werden, da aus den
verschiedenen vorhandenen Rechnungen nicht ersichtlich ist, welche Methode
\textsc{Gauss} schliesslich als zweckmässigste angesehen haben mochte. Die ausführliche
Hülfstafel, die Seite 560 erwähnt ist, steht im Handbuch Bc; ihr Abdruck erschien
unnöthig. Die numerischen Werthe zur Vergleichung mit den Oppositionen Seite 560
stehen ebenfalls verstreut im Handbuch Bc; die Störungsbeträge sind hier nach den
oben erwähnten Manipeln berechnet, wobei \textsc{Gauss} von \textsc{Wachter} unterstützt
wurde. Es finden sich aber nur die Rechnungen für die vier S. 560 behandelten Elemente vor.

\par\noindent\hspace*{\parindent}Die Notiz [12.], Seite 561 infolge eines Versehens mit [11.] bezeichnet,
enthält die Vergleichung mit den ersten neun Oppositionen, die \textsc{Gauss} in
seinem Brief an \textsc{Bessel} vom 21. März 1843 (S. 438) erwähnt; auch sie ist auf
Blättern ausgeführt, denen nur die Zahlen entnommen werden konnten; auf diesen
Blättern hat \textsc{Gauss} später noch einige Zahlen hinzugefügt, die sich auf
die zehnte Opposition beziehen. Ausser
\SourceLeaf{615}diesen hier abgedruckten finden sich, wie bereits Eingangs dieser Bemerkungen
erwähnt, noch sehr viele weitere Vergleichungen, auch mit spätern Oppositionen.
\par
\endgroup
\clearpage
\endgroup
